% s-318-p318-m003; original PDF 318, 319
\Aphorism{s-318-p318-m003}{374}
Dass Imperative in Fleisch und Blut übergehen, dass befohlene Acte in Willensacte und diese in Reflexacte umgeschmolzen werden, dass Autorität an Einsicht und diese an das „Unbewusste“ den Oberbefehl abtrete: dies ist die natürliche Entwicklung, das natürliche Ziel jeder Ethik, die aus dem Leben herauswächst und demgemäss ihre Selbstverneinung will. Kant sucht hingegen mit dem ganzen Fanatismus des Naturverbesserers einen Stillstand in jener Entwicklung zu erzwingen und die Ethik auf der niedrigsten Vorstufe festzunageln: die „Pflicht“ soll Pflicht bleiben und sich beileibe nicht zur Neigung oder zur reflexmässigen Unwillkür degradiren. Soll! als ob der Mensch nichts anderes auf der Welt zu thun hätte als zu sollen! ganz abgesehen davon, \Sperr{was }er eigentlich nach Kant soll. Aber selbst angenommen, die Moral gebiete sein Bestes, so thäte er doch wohl daran, dieses Gebot nicht ewig als Gebot und „fixe Idee“ über sich hängen zu lassen, sondern seiner bewussten und schliesslich physiologischen Intelligenz einzuverleiben, um Bewusstsein und Idealismus für vielleicht noch höhere Aufgaben freizubekommen.
\par

